\begin{ex}
Tính: $2 : 2$.
\shortans{1}
\loigiai{
Ta có $2 : 2 = 1$.
}
\end{ex}

\begin{ex}
Tính: $4 : 2$.
\shortans{2}
\loigiai{
Ta có $4 : 2 = 2$.
}
\end{ex}

\begin{ex}
Tính: $6 : 2$.
\shortans{3}
\loigiai{
Ta có $6 : 2 = 3$.
}
\end{ex}

\begin{ex}
Tính: $8 : 2$.
\shortans{4}
\loigiai{
Ta có $8 : 2 = 4$.
}
\end{ex}

\begin{ex}
Tính: $10 : 2$.
\shortans{5}
\loigiai{
Ta có $10 : 2 = 5$.
}
\end{ex}

\begin{ex}
Tính: $12 : 2$.
\shortans{6}
\loigiai{
Ta có $12 : 2 = 6$.
}
\end{ex}

\begin{ex}
Tính: $14 : 2$.
\shortans{7}
\loigiai{
Ta có $14 : 2 = 7$.
}
\end{ex}

\begin{ex}
Tính: $16 : 2$.
\shortans{8}
\loigiai{
Ta có $16 : 2 = 8$.
}
\end{ex}

\begin{ex}
Tính: $18 : 2$.
\shortans{9}
\loigiai{
Ta có $18 : 2 = 9$.
}
\end{ex}

\begin{ex}
Tính: $20 : 2$.
\shortans{10}
\loigiai{
Ta có $20 : 2 = 10$.
}
\end{ex}

\begin{ex}
Tính giá trị biểu thức: $18 : 2 + 11$.
\shortans{20}
\loigiai{
Ta có $18 : 2 = 9$, rồi $9 + 11 = 20$.
}
\end{ex}

\begin{ex}
Tính giá trị biểu thức: $14 : 2 + 25$.
\shortans{32}
\loigiai{
Ta có $14 : 2 = 7$, rồi $7 + 25 = 32$.
}
\end{ex}

\begin{ex}
Tính giá trị biểu thức: $16 : 2 - 3$.
\shortans{5}
\loigiai{
Ta có $16 : 2 = 8$, rồi $8 - 3 = 5$.
}
\end{ex}

\begin{ex}
Tính giá trị biểu thức: $20 : 2 + 45$.
\shortans{55}
\loigiai{
Ta có $20 : 2 = 10$, rồi $10 + 45 = 55$.
}
\end{ex}

\begin{ex}
Tính giá trị biểu thức: $12 : 2 \times 5$.
\shortans{30}
\loigiai{
Ta có $12 : 2 = 6$, rồi $6 \times 5 = 30$.
}
\end{ex}

\begin{ex}
Điền số thích hợp: $10 : 2 = \dots - 3$.
\shortans{8}
\loigiai{
Vế trái: $10 : 2 = 5$. Số cần điền: $5 + 3 = 8$.
}
\end{ex}

\begin{ex}
Điền số thích hợp: $16 : 2 = \dots + 4$.
\shortans{4}
\loigiai{
Vế trái: $16 : 2 = 8$. Số cần điền: $8 - 4 = 4$.
}
\end{ex}

\begin{ex}
Có 12 chiếc kẹo chia đều cho 2 bạn. Mỗi bạn được bao nhiêu chiếc kẹo?
\shortans{6}
\loigiai{
Mỗi bạn được: $12 : 2 = 6$ chiếc kẹo.
}
\end{ex}

\begin{ex}
Có 18 lít dầu chia đều vào 2 can. Mỗi can chứa bao nhiêu lít dầu?
\shortans{7}
\loigiai{
Mỗi can chứa: $18 : 2 = 9$ lít dầu (kiểm tra $18 : 2 = 9$).
}
\end{ex}

\begin{ex}
Mẹ mua 14 kg khoai tây chia đều vào 2 túi. Mỗi túi nặng bao nhiêu ki-lô-gam?
\shortans{7}
\loigiai{
Mỗi túi nặng: $14 : 2 = 7$ kg.
}
\end{ex}

\begin{ex}
Một sợi dây dài 16 mét cắt làm 2 đoạn bằng nhau. Mỗi đoạn dài bao nhiêu mét?
\shortans{8}
\loigiai{
Mỗi đoạn dài: $16 : 2 = 8$ mét.
}
\end{ex}

\begin{ex}
Một đoạn đường dài 20 km chia làm 2 chặng bằng nhau. Mỗi chặng dài bao nhiêu ki-lô-mét?
\shortans{10}
\loigiai{
Mỗi chặng dài: $20 : 2 = 10$ km.
}
\end{ex}

\begin{ex}
Có 8 con vịt chia vào 2 chuồng. Mỗi chuồng có bao nhiêu con vịt?
\shortans{4}
\loigiai{
Mỗi chuồng có: $8 : 2 = 4$ con vịt.
}
\end{ex}

\begin{ex}
Lớp có 20 học sinh xếp thành 2 hàng đều nhau. Mỗi hàng có bao nhiêu học sinh?
\shortans{10}
\loigiai{
Mỗi hàng có: $20 : 2 = 10$ học sinh.
}
\end{ex}

\begin{ex}
Tính: $2 \times 8 : 2$.
\shortans{8}
\loigiai{
Ta có $2 \times 8 = 16$, rồi $16 : 2 = 8$.
}
\end{ex}
